\section{Results}
\label{sec:results}

\subsection{Average response of house prices}

\begin{figure}[htbp]
  \centering
  \includegraphics[width=0.85\textwidth]{Fig01_irf_baseline.pdf}
  \caption{Average response of real house prices to a one basis point
  tightening surprise, in percent, with a 95 percent confidence band.
  Baseline specification with dynamic controls, country fixed effects
  and Driscoll Kraay standard errors, eleven euro area countries,
  1999Q4 to 2025Q4.}
  \label{fig:baseline}
\end{figure}

Figure \ref{fig:baseline} shows the estimated response of real house prices to a
one basis point tightening surprise, averaged across the eleven countries. The
response is close to zero at every horizon. It is mildly negative over the first
five quarters, with a trough of about $-0.03$ percent around one year, and then
drifts back towards zero and turns weakly positive at longer horizons. None of
the coefficients is statistically different from zero. The confidence bands are
wide and widen with the horizon, which is what one expects when the identifying
variation comes from a single euro area wide shock observed in 108 quarters.

To put the magnitude in perspective, the standard deviation of the quarterly
shock is about 4.6 basis points (Table \ref{tab:summary}). A one standard
deviation tightening surprise therefore moves house prices by roughly a tenth of
a percentage point in the point estimates, which is economically small as well
as statistically indistinguishable from zero.

\subsection{Does mortgage market structure matter?}

The central question is whether the response differs between countries where
mortgages are mostly variable rate and countries where they are mostly fixed
rate. Table \ref{tab:main} reports the interaction between the shock and the
adjustable rate indicator at horizons of zero, four and eight quarters, and
Figure \ref{fig:interaction} plots it over the full twelve quarter horizon.

{\footnotesize\input{tables/Table02_lp_main}}

\begin{figure}[htbp]
  \centering
  \includegraphics[width=0.85\textwidth]{Fig02_irf_interaction.pdf}
  \caption{Differential response of variable rate countries, in
  percentage points per basis point of tightening surprise. The solid
  line is the baseline specification with dynamic controls, with a 95
  percent confidence band. The dotted line is the same specification
  without the dynamic controls. Country fixed effects and Driscoll
  Kraay standard errors throughout.}
  \label{fig:interaction}
\end{figure}

In the baseline specification the interaction is positive at every horizon,
rising from about $0.02$ on impact to $0.06$ after twelve quarters. A positive
coefficient means that house prices in variable rate countries fall by
\emph{less} than in fixed rate countries after a tightening surprise, which is
the opposite of what the budget channel predicts. The coefficient is marginally
significant at the ten percent level on impact, where the ratio of estimate to
standard error reaches about $1.9$, and is insignificant at every other
horizon. None of the shock or interaction coefficients reaches the
conventional five percent threshold. The
estimates therefore cannot distinguish the hypothesised amplification from no
differential effect at all, nor from a modest effect in the opposite direction.

The dotted line in Figure \ref{fig:interaction} shows the same interaction when
the dynamic controls are removed. It is negative at every horizon and reaches
about $-0.12$ after nine quarters, which is the sign the hypothesis predicts.
This coefficient is also statistically insignificant, with a largest ratio of
estimate to standard error of about $1.1$. The two specifications therefore
disagree on the sign of the differential effect while agreeing that it cannot be
distinguished from zero.

I report both specifications rather than choosing between them after seeing the
results. The specification with dynamic controls is the baseline because
including lags of the shock and of the dependent variable is standard practice
in local projections, and that choice was made on methodological grounds before
the estimates were compared. The sensitivity of the sign to this choice is
itself informative. It shows that the differential effect is not identified
precisely enough for the data to pin down even its direction.

\subsection{Interpretation}

Taken together, the results provide no robust evidence that the structure of the
national mortgage market amplifies the transmission of ECB monetary policy to
house prices. The point estimates are small, none of them reaches significance
at the five percent level in either specification, and their sign is not
stable across reasonable modelling choices. This is a null result rather than
evidence against the mechanism. The confidence intervals are wide enough to
contain both the amplification predicted by theory and effects of the opposite
sign, so the exercise mainly shows that a common euro area shock and eleven
national mortgage systems identify very little.

Two further findings are reported in the appendix. First, the information
correction applied to the shock matters a great deal. Using the raw interest
rate surprise instead of the information cleaned series produces a positive
house price response to tightening that grows with the horizon, which is not
consistent with any standard transmission mechanism and is a typical sign
of central bank information effects contaminating the surprise. Second, splitting
the sample at 2009 shows that the contemporaneous average response is negative
and statistically significant in the later part of the sample, which suggests
that some of the weakness in the full sample estimates comes from averaging over
periods with very different monetary regimes.
